% ===================================================================
%  TABEL Nr. 7 — Küla talvel. Talumaja kevadel.
%  Traduction 2026 d'apres texto/io, controlee sur texto/fr et
%  texto/fi. Consigne : voir texto/et/10-tabel-01.tex.
%
%  DEUX MOTS DE L'ALINEA 10 SE TIENNENT COTE A COTE ET NE VIENNENT
%  PAS DU MEME MONDE : LE MOULIN ET LE MEUNIER.
%
%    « veski »  le moulin, mot estonien, tire de « vesi » l'eau ;
%    « mölder » le meunier, l'allemand Müller pris tel quel.
%
%  LA MACHINE A GARDE SON NOM ; L'HOMME QUI LA FAIT TOURNER N'EN A
%  PAS GARDE. Ce n'est pas un hasard de vocabulaire, c'est un fait de
%  droit : dans les provinces baltes le moulin appartenait au manoir,
%  la mouture etait un monopole seigneurial, et le meunier etait
%  l'homme du seigneur, nomme par lui et responsable devant lui.
%
%  ET LA COMPARAISON EST SUR LA MEME PAGE. Le forgeron de l'alinea 5
%  s'appelle « sepp », mot estonien du fonds le plus ancien, celui-la
%  meme que le tableau 5 a montre servant a fabriquer six noms de
%  metier. Pourquoi le forgeron a-t-il garde le sien et le meunier
%  pas ? PARCE QUE PERSONNE NE POSSEDAIT LA FORGE. Le village
%  entretenait son forgeron ; le moulin, lui, avait un proprietaire.
%
%  LE TABLEAU 11 ROMANCHE AVAIT TROUVE QUE QUI FONDE NOMME, ET LE 12
%  L'AVAIT VERIFIE SUR LE CHEMIN DE FER. L'estonien ajoute le cas que
%  la Suisse ne pouvait pas fournir, n'ayant pas connu le servage :
%  QUI POSSEDE NOMME AUSSI, ET IL NOMME L'HOMME AVANT LA CHOSE. Un
%  proprietaire ne rebaptise pas la meule, il rebaptise celui qui
%  l'actionne, parce que c'est celui-la qui lui doit des comptes.
%
%  La suite de la page le confirme sans qu'on la force : le
%  « kingsepp » le cordonnier, le « vankrisepp » le charron, le
%  « sadulsepp » le bourrelier sont tous des composes en « -sepp » —
%  ce sont les metiers qu'un village paie lui-meme.
%
%  TROIS INVERSIONS SUR CETTE SEULE PAGE, ET C'EST TROIS FOIS LA MEME
%  CAUSE : la chaine de genitifs estonienne, deja rencontree aux
%  tableaux 3 et 5. « oja jää » la glace du ruisseau, « mära
%  allapanu » la litiere de la jument, « koera kuudi » la niche du
%  chien — le possesseur passe devant et l'ido range l'inverse. Le
%  remede est le meme les trois fois et il est desormais mecanique :
%  on garde le mot possede en tete et l'on rejette le possesseur dans
%  une relative — « la glace, qui couvrait le ruisseau », « la
%  litiere, qui appartient a la jument », « la niche, ou vit le
%  chien ». Le genitif estonien est la seule difficulte d'ordre de
%  toute la colonne, et elle a une solution unique.
% ===================================================================

%%K t07-noto-1 noto
\VUnotes{143.2mm}{%
(*) Söögi üksikasjade kohta vaata tabeleid 6 ja 13.%
}

%%K t07-apar-1 apar
\VUsaut{4.0mm}
\VUcentre{13.2pt}{90}{TEINE SEERIA}
\VUsaut{3.0mm}
\VUfilet{16mm}
\VUsaut{5.6mm}
\VUtitre{15.0pt}{60}{TABEL Nr. 7}
\VUsaut{3.0mm}

%%K t07-tit-1 sub
\VUcentre{12.0pt}{0}{\textit{Esimene stseen.}}
\VUsaut{3.0mm}

%%K t07-tit-2 sub
\VUpk{10.2pt}{40}{Küla talvel.}
\VUsaut{4.0mm}

%%K t07-c1-01-1 p
1. --- Uusaastavaheajal saatsin ma oma isa Uuslinna, väikesesse külla, kus tal oli
vaja ajada mõned kaubaasjad.

%%K t07-c1-02-1 p
2. --- Me kasutasime \VUgras{postitõlda} \textsuperscript{(11)}, mis käib linna ja selle
aleviku vahet; kuid kuna oli väga külm, ei olnud see sõit vähe mugavas sõidukis eriti
meeldiv.

%%K t07-c1-03-1 p
3. --- Seepärast tundsime me tõelist rõõmu, kui nägime \VUgras{kutsarit} oma vankrit
platsil peatamas, \textit{Valge Karu} \VUgras{võõrastemaja} \textsuperscript{(2)} ees.
\VUgras{Võõrastemaja peremees} \textsuperscript{(4)} ootas meid oma maja lävel,
rahulikult oma \VUgras{piipu} suitsetades.

%%K t07-c1-03-2 p
Ta kutsus meid sisse ja ma ei lasknud end paluda, et end sisse seada raagude tule ees,
mis koldes praksus. Siis naersin ma väga terava \VUgras{põhjatuule} üle, mis pani vana
\VUgras{sildi} \textsuperscript{(3)} kriuksuma ja kandis mustades keeristes ära
\VUgras{suitsu} \textsuperscript{(1)}, mis korstnast väljus.

%%K t07-c1-04-1 p
4. --- Oli lõunasöögi aeg. Kauem ootamata sõime me hästi lihtsa, kuid maitsva söögi,
mida meile pakuti (*).

%%K t07-tit-3 sub
\VUpk{10.2pt}{40}{Mõned külaskäigud.}
\VUsaut{3.0mm}

%%K t07-c1-05-1 p
5. --- Pärast lõunasööki saatsin ma oma isa, kes on kindlustusagent, tema klientide
juurde. Kõigepealt külastasime me \VUgras{seppa} \textsuperscript{(5)}, kes elab
võõrastemaja lähedal. Ta oli ametis hobuse rautamisega. Ma imetlesin tema abilise
tugevust, kes hoidis kindlalt oma nahkpõlle peal hobuse \VUgras{kabja}, sellal kui
sepp kinnitas \VUgras{raua} \textsuperscript{(6)} tugevate haamrilöökidega.

%%K t07-c1-06-1 p
6. --- Siis läksime me küla \VUgras{kingsepa} \textsuperscript{(7)}, vana Jaakobi
juurde. Kuigi tal on sildiks väga ilus \VUgras{saabas}, leppis ta sellega, et pani uued
tallad vanale paarile auklikele kingadele.

%%K t07-c1-06-2 p
Vanamees ütles meile naerdes: «Linnas olin ma «kingavalmistaja» ja mul ei olnud
kliente. Siin olen ma «kingaparandaja» ja tööd on palju.»

%%K t07-c1-07-1 p
7. --- Ta rääkis meile oma naabrist \VUgras{habemeajajast} \textsuperscript{(8)}, kelle
klientuur ei ole rahul. Tema \VUgras{habemenoad} on alati täkilised ja tema
\VUgras{käärid} ei lõika kunagi hästi.

Kuid kuna ta on küla ainus habemeajaja, ollakse muidugi sunnitud laskma tal end ajada
ja juukseid lõigata. Pealegi on ta ihne ja ebasõbralik mees.

%%K t07-c1-08-1 p
8. --- Õnneks ei sarnane teised \VUgras{naabrid} temaga. Näiteks \VUgras{vankrisepp}
\textsuperscript{(9)} on heasüdamlik mees, töökas ja abivalmis. On rõõm lasta tal
vankrit parandada. Tema töö on alati õigeks ajaks valmis.

%%K t07-c1-09-1 p
9. --- Vana Jaakob ei kurtnud ka \VUgras{sadulsepa} \textsuperscript{(10)} üle, keda ta
mõnikord aitab \VUgras{rakmeid} õmmelda, kui tööga on kiire.

%%K t07-c1-09-2 p
Mu isa küsis temalt tema perekonna kohta: «Kõigil läheb hästi. Mu lapselaps on
\VUgras{koolis} \textsuperscript{(12)}. Tema \VUgras{õpetajanna} \textsuperscript{(13)} on
temaga väga rahul, ta on hea \VUgras{õpilane} \textsuperscript{(14)}. Ta ei sarnane mu
üleannetu poja moodi; see mõtleb ainult mängimisele. \VUgras{Õpetaja}
\textsuperscript{(15)} ei suuda tal midagi selgeks teha.»

%%K t07-tit-4 sub
\VUsaut{3.0mm}
\VUpk{10.2pt}{40}{Külajutud.}
\VUsaut{3.0mm}

%%K t07-c1-10-1 p
10. --- Vanamees jutustas meile siis küla uudised. Hakatakse ehitama uut kooli, sest
see, mis on \VUgras{vallamajas}, on liiga väikeseks jäänud. Pealegi on see lihtne, sest
piisab osta osa \VUgras{möldri} aiast. See tuleb vaesele mehele väga kasuks. Külma
pärast ei suuda \VUgras{veski} \textsuperscript{(16)} \VUgras{veskikivid} enam nisu
jahvatada, sest \VUgras{ratta} \textsuperscript{(17)} on kinni pannud \VUgras{jää}
\textsuperscript{(25)}, mis kattis \VUgras{oja} \textsuperscript{(23)}.

%%K t07-c1-11-1 p
11. --- «Ehitistest rääkides, kas te nägite meie uut \VUgras{kirikut}
\textsuperscript{(18)}? See on väga maaliline oma lumemantli all. \VUgras{Vareseid}
\textsuperscript{(19)}, kes ei tunne enam ära oma vana kellatorni, istuvad kurvalt
\VUgras{surnuaia} \textsuperscript{(20)} suure puu otsas.»

%%K t07-c1-12-1 p
12. --- See mõte surnuaiast tuletas kingsepale meelde inimesed, keda mu isa oli tundnud
ja kes eelmisel aastal surid. «Kui palju \VUgras{haudu} \textsuperscript{(21)}, mu hea
härra, on teiste juurde lisandunud \VUgras{küpressi} \textsuperscript{(22)} all! Surm
niitis meie vana notari. Kõik külaelanikud viibisid tema \VUgras{matusel}.»

%%K t07-c1-13-1 p
13. --- «Vaat et tema järeltulija uisutab ojal. Ta tuli just laskma mul parandada oma
\VUgras{uisu} \textsuperscript{(26)} rihma, mis oli katki.»

%%K t07-c1-14-1 p
14. --- «Vaat et ka oja kaldal on uus \VUgras{külavaht} \textsuperscript{(27)}. Ta on
endine sandarm, kes hirmutab küla üleannetuid poisse. Nad põgenevad kohe, kui nad
näevad tema \VUgras{sääriseid} \textsuperscript{(29)} ja tema \VUgras{mundrimütsi}
\textsuperscript{(28)}.»

%%K t07-c1-15-1 p
15. --- «Ah! vaat et vana Joosep veab \VUgras{koormat küttepuid} \textsuperscript{(30)}
pagari jaoks. --- See on tõsi, ütles mu isa, ma tunnen ära tema \VUgras{muulade}
\textsuperscript{(31)} rakendi. Mul on tõesti vaja temaga rääkida, ma kasutan juhust ära.
Nägemiseni, isa Jaakob.»

%%K t07-c1-16-1 p
16. --- Just kui me välja tulime, jätsid küla poisid kooli ja tormasid joostes
\VUgras{liumäele} \textsuperscript{(33)}, riskides \VUgras{kukkumisel}
\textsuperscript{(34)} maha kukkuda ja jäset murda. Üks neist võttis vankrisepa juurest
oma \VUgras{kelgu} \textsuperscript{(32)}, pani sellele istuma ühe oma kaaslastest ja
läks joostes minema.

%%K t07-c1-17-1 p
17. --- Teised tormasid \VUgras{lumehange} \textsuperscript{(37)} \VUgras{silla}
\textsuperscript{(40)} juures ja hakkasid pommitama \VUgras{lumememme}
\textsuperscript{(39)}, kelle nad olid eelmisel päeval püstitanud ja kellele nad olid
pannud kübara, piibu ja üle õla riputatud koolikoti.

%%K t07-c1-18-1 p
18. --- Nad ei hooli külmetavast tuulest ega pakasest. Enamikul ei ole isegi
\VUgras{salli} \textsuperscript{(35)}, ja et pärast \VUgras{lumepallide}
\textsuperscript{(38)} lahingut end soojendada, piisab neile peotäiest väga kuumadest
\VUgras{kastanitest} \textsuperscript{(42)}, mis on ostetud vanalt \VUgras{müüjannalt}
Jaagupilt, kes väriseb külmast oma liiga õhukese \VUgras{räbala rätiku}
\textsuperscript{(43)} all.

%%K t07-tit-5 sub
\VUsaut{3.0mm}
\VUcentre{12.0pt}{0}{\textit{Teine stseen.}}
\VUsaut{3.0mm}

%%K t07-tit-6 sub
\VUpk{10.2pt}{40}{Talumaja kevadel.}
\VUsaut{4.0mm}

%%K t07-c2-01-1 p
1. --- Vaatamata kõigile talve rõõmudele tervitame me sügava rõõmuga \VUgras{kevade}
tagasitulekut.

%%K t07-c2-01-2 p
Kui rõõmustava pildi pakub talu õu õhtul maikuus!

%%K t07-c2-02-1 p
2. --- Silmapiiril laskub \VUgras{päike} \textsuperscript{(1)} aeglaselt, uputades
\VUgras{maastiku} \textsuperscript{(3)} oma purpursetesse \VUgras{kiirtesse}
\textsuperscript{(2)}. Usin \VUgras{kündja} \textsuperscript{(6)} kiirustab oma
\VUgras{põldu} \textsuperscript{(4)} kündma.

%%K t07-c2-02-2 p
Oma \VUgras{adraga} \textsuperscript{(5)}, mis on rakendatud tugeva \VUgras{hobuse}
\textsuperscript{(7)} ette, veab ta \VUgras{vao} \textsuperscript{(8)}, millesse
\VUgras{külvaja} \textsuperscript{(9)} viskab täie käega \VUgras{seemne}, mis annab
kuldsed viljapead.

%%K t07-c2-03-1 p
3. --- \VUgras{Niitjad} \textsuperscript{(11)}, varustatud oma vikatitega, on niitnud
heinamaa rohu, heinalised on \VUgras{heina} \textsuperscript{(10)} kuivatanud,
keerates seda \VUgras{hangudega} \textsuperscript{(12)} ja rehadega, ja vaat et nad
tulevadki tagasi.

%%K t07-c2-03-2 p
Ühed kõnnivad härgade veetava raske vankri ees; nad kannavad oma tööriista õlal.
Teised, laisemad, on end mugavalt seadnud lõhnavale heinale.

%%K t07-c2-04-1 p
4. --- Möödudes tervitavad nad sõbralikult \VUgras{aednikku} \textsuperscript{(16)}, kes
on ametis \VUgras{viljapuuaias} \textsuperscript{(13)} \VUgras{viljapuude} lõikamise ja
\VUgras{pirnipuude} \textsuperscript{(14)} pookimisega. \VUgras{Ploomipuud}
\textsuperscript{(15)} hakkavad õitsema, ja hästi sõnnikustatud \VUgras{juurviljaaias}
\textsuperscript{(17)} on juurviljad, \VUgras{kapsad} ja salatid juba heas seisus.

%%K t07-c2-04-2 p
Väiksel müüril laotab ilus \VUgras{paabulind} \textsuperscript{(18)} oma \VUgras{saba},
et lasta imetleda selle imeilusaid \VUgras{sulgi}.

%%K t07-tit-7 sub
\VUpk{10.2pt}{40}{Talu õu.}
\VUsaut{3.0mm}

%%K t07-c2-05-1 p
5. --- \VUgras{Talli} \textsuperscript{(19)} uksed on avatud ja näha on sõim ning
\VUgras{allapanu} \textsuperscript{(23)}, mis kuulub \VUgras{märale}
\textsuperscript{(21)}, kelle \VUgras{tallipoiss} \textsuperscript{(20)} on äsja lahti
rakendanud. \VUgras{Varss}
\textsuperscript{(22)}, nii naljakas oma pikkade jalgadega, järgneb hüpeldes oma emale.

%%K t07-c2-06-1 p
6. --- \VUgras{Kaevu} \textsuperscript{(29)} juures lähevad \VUgras{lehmad}
\textsuperscript{(25)} jooma puust \VUgras{künasse} \textsuperscript{(26)}, mis on neile
jooginõuks. \VUgras{Vasikas} \textsuperscript{(24)} kasutab juhust imemiseks, ja
\VUgras{pull} \textsuperscript{(27)}, kes tunneb heina head lõhna, ammub rõõmsalt ja
tõstab uhkelt pead oma võimsate \VUgras{sarvedega} \textsuperscript{(28)}.

%%K t07-c2-07-1 p
7. --- Kõik töötavad. \VUgras{Teenijatüdruk} \textsuperscript{(32)} hoiab käes kaevu
\VUgras{köit} \textsuperscript{(31)}. Ta ammutab vett \VUgras{ämbriga}
\textsuperscript{(30)}, et karja joota. \VUgras{Küüni} \textsuperscript{(33)} juures
riisub mees õlgi, ja \VUgras{elumaja} \textsuperscript{(34)} ukse ees ketrab vana
\VUgras{ketraja} \textsuperscript{(35)} oma voki ja oma \VUgras{kedervarrega}
\VUgras{lina} või \VUgras{kanepit} nagu vanasti.

%%K t07-tit-8 sub
\VUsaut{3.0mm}
\VUpk{10.2pt}{40}{Talunik.}
\VUsaut{3.0mm}

%%K t07-c2-08-1 p
8. --- \VUgras{Talunik} \textsuperscript{(36)} juhib kõiki neid töid ja annab juhiseid
oma sulastele. Ta soovitab katuse alla panna \VUgras{äke} \textsuperscript{(40)} ja
\VUgras{kirka} \textsuperscript{(39)}, mis unustati \VUgras{kuudi}
\textsuperscript{(38)} juurde, kus elab \VUgras{koer} \textsuperscript{(37)}.

%%K t07-c2-09-1 p
9. --- Tema hüüded hirmutavad \VUgras{tuvi} \textsuperscript{(41)}, kes lendab kõigest
jõust \VUgras{tuvipuuri} \textsuperscript{(42)}, ja \VUgras{kanad}
\textsuperscript{(43)}, kes kiirustavad tagasi \VUgras{kanalasse} \textsuperscript{(44)}.

%%K t07-c2-10-1 p
10. --- \VUgras{Lammaslauda} \textsuperscript{(48)} ees on vana \VUgras{karjus}
\textsuperscript{(45)}, kes toob tagasi oma \VUgras{lambakarja} \textsuperscript{(49)}.
Hoides oma \VUgras{karjasekeppi} \textsuperscript{(46)}, mis tal kunagi ei puudu, nagu ka
oma pleekinud \VUgras{kuube} \textsuperscript{(47)}, viib ta lauta \VUgras{oinad}
\textsuperscript{(50)}, \VUgras{utid} \textsuperscript{(53)}, \VUgras{talled}
\textsuperscript{(54)}, \VUgras{kitsed} \textsuperscript{(51)} ja \VUgras{kitsetalled}
\textsuperscript{(52)}. Tema truu \VUgras{koer} \textsuperscript{(55)} Argus tuleb
viimasena ja kiirustab oma haukumisega mahajääjaid.

%%K t07-c2-11-1 p
11. --- Kogu see lärm ja liikumine ei häiri hea \VUgras{lüpsilehma}
\textsuperscript{(56)} rahu, kes kõlistab oma \VUgras{kella} \textsuperscript{(57)}, sellal
kui teda lüpstakse \VUgras{lauda} \textsuperscript{(58)} ukse ees.

%%K t07-c2-12-1 p
12. Ta näib mõistvat, et ta peab andma oma piima, et \VUgras{taluperenaine}
\textsuperscript{(60)} saaks valmistada \VUgras{võid} \VUgras{võimasinas}
\textsuperscript{(61)} või suurepäraseid \VUgras{juuste} \textsuperscript{(64)}, mis
tilguvad \VUgras{tõrre} \textsuperscript{(63)} peale pandud laualt.

%%K t07-c2-13-1 p
13. --- Kogu \VUgras{piimakambrit} \textsuperscript{(59)} hoiab väga puhtana
\VUgras{talusulane} \textsuperscript{(65)}, kes on äsja pesnud \VUgras{piimanõud}
\textsuperscript{(62)} suures ämbris, mida ta käes kannab.

%%K t07-tit-9 sub
\VUsaut{3.0mm}
\VUpk{10.2pt}{40}{Kodulinnud.}
\VUsaut{3.0mm}

%%K t07-c2-14-1 p
14. --- Oma paksude puukingadega kõnnib ta pisut kohmakalt ja ajab nii põgenema
\VUgras{kalkuni} \textsuperscript{(68)}, \VUgras{pardid} \textsuperscript{(66)},
\VUgras{isapardid} \textsuperscript{(67)} ja \VUgras{haned} \textsuperscript{(69)}, kes
avavad laialt \VUgras{oma noka} \textsuperscript{(70)} ja rahutult kisavad.

%%K t07-c2-14-2 p
\VUgras{Kukk} \textsuperscript{(71)}, võitlushimulisem, ajab püsti oma punase
\VUgras{hari} \textsuperscript{(72)}, raputab oma \VUgras{saba} \textsuperscript{(73)}
sulgi ja laseb kuulda kõlava «kikerikii».

%%K t07-c2-15-1 p
15. --- \VUgras{Kanaema} \textsuperscript{(74)} otsib \VUgras{sõnnikul}
\textsuperscript{(81)} koos oma \VUgras{tibudega} \textsuperscript{(75)}.
\VUgras{Põrsas} \textsuperscript{(78)} põgeneb oma ema, tohutu \VUgras{emise}
\textsuperscript{(77)} juurde, keda me näeme sealauda kõrval.

%%K t07-c2-16-1 p
16. --- Oma sulgudes söövad \VUgras{küülikud} \textsuperscript{(79)} ahnelt õrna
\VUgras{rohtu} \textsuperscript{(80)}, mida neile toob taluniku tütar. Jaanika armastab
väga oma küülikuid ja tuleb iga päev, olles kanalast värsked \VUgras{munad}
\textsuperscript{(76)} toonud, oma väikesi sõpru vaatama.

\VUsaut{6.0mm}
\VUfilet{22mm}
